\section{Fundamentação}

\subsection{Estrutura eletrônica em computadores quânticos}

Na aproximação de Born--Oppenheimer, os núcleos são tratados como fixos para
cada geometria e resolve-se o problema eletrônico. Em segunda quantização, a
Hamiltoniana é escrita em termos de operadores fermiônicos de criação e
aniquilação. Uma transformação como Jordan--Wigner os converte em uma soma de
produtos de Pauli,
\begin{equation}
 H_q = \sum_j c_j P_j, \qquad
 P_j \in \{I,X,Y,Z\}^{\otimes n},
\end{equation}
que pode ser medida por um computador quântico. O número de qubits depende do
número de spin-orbitais mantidos e o número de termos $P_j$ oferece um descritor
simples do tamanho da Hamiltoniana.

As energias exatas no espaço orbital escolhido são obtidas por \textit{Full
Configuration Interaction} (FCI). Quando se emprega um espaço ativo, CASCI
diagonaliza a Hamiltoniana somente nos orbitais selecionados. Neste estudo,
essas energias são referências para calcular o erro do VQE, e não entradas dos
modelos de aprendizado. Os cálculos clássicos foram realizados com PySCF
\cite{sun2018pyscf}.

\subsection{Variational Quantum Eigensolver}

O VQE explora o princípio variacional: para qualquer estado normalizado, o
valor esperado da Hamiltoniana limita superiormente a energia do estado
fundamental. O algoritmo prepara um estado parametrizado
$\lvert\psi(\boldsymbol{\theta})\rangle$ e minimiza
\begin{equation}
 E(\boldsymbol{\theta}) =
 \langle\psi(\boldsymbol{\theta})\rvert H
 \lvert\psi(\boldsymbol{\theta})\rangle .
\end{equation}
O processador quântico estima os valores esperados dos termos de Pauli,
enquanto um otimizador clássico atualiza $\boldsymbol{\theta}$
\cite{mcclean2016theory}. Essa arquitetura híbrida foi demonstrada em moléculas
pequenas com circuitos eficientes em hardware \cite{kandala2017hardware}.

Ansätze fisicamente inspirados, como UCCSD e sua versão pareada PUCCSD,
incorporam excitações da função de onda eletrônica. Em contrapartida,
RealAmplitudes e EfficientSU2 são circuitos heurísticos construídos com camadas
de rotações e emaranhamento. A profundidade, controlada aqui por
\textit{repetitions} (\textit{reps}), aumenta a expressividade, mas também o
número de parâmetros e o custo do circuito. Circuitos muito expressivos podem
apresentar regiões de gradientes muito pequenos, conhecidas como
\textit{barren plateaus} \cite{mcclean2018barren}.

COBYLA e SLSQP realizam buscas determinísticas com aproximações locais,
enquanto SPSA estima a direção de atualização por perturbações estocásticas.
Logo, a configuração do VQE envolve compromissos entre estrutura física,
profundidade e comportamento do otimizador. Não existe uma escolha
universalmente superior para todas as moléculas e geometrias.

Neste estudo, uma configuração atinge precisão química quando
$\lvert E_{\mathrm{VQE}}-E_{\mathrm{ref}}\rvert\leq 1$ kcal/mol. Esse limiar
transforma a qualidade energética em um critério prático de sucesso, mas o
modelo é treinado para regressão do erro contínuo, preservando a ordenação
entre candidatos \cite{petruzielo2012chemical}.

\subsection{Seleção supervisionada de configurações}

O problema de seleção de algoritmos consiste em escolher, a partir de
características de uma instância, o método com melhor desempenho esperado
\cite{rice1976algorithm}. Aqui, cada instância é uma geometria molecular e cada
candidato é uma configuração VQE. Em vez de classificar diretamente um único
vencedor, estima-se uma função de custo para todos os candidatos e utiliza-se a
ordenação prevista.

Modelos lineares oferecem um baseline interpretável; Random Forest representa
relações não lineares em dados tabulares \cite{breiman2001random}; e redes
multicamadas aproximam mapeamentos não lineares por camadas sucessivas
\cite{hornik1989multilayer}. A regressão fornece uma pontuação comparável para
todos os candidatos. Ordená-los por erro previsto transforma o problema em
recomendação top-$k$: executam-se poucos candidatos e conserva-se o melhor
resultado observado entre eles.
